% Propietario conservado: /Users/ruben/Documents/New project/output/REVISION_ARGUMENTAL_APERTURAS_CIERRES_20260915/VII/delivery/MANUSCRIPT_VII_ES_EN_COMPLETE/source_es/manuscrito/31_corriente_y_respuesta_cartan.tex
% Artículo VII. Fragmento de cuerpo: corriente y respuesta Cartan--Holst.
% RESULTADO_RECUPERADO: acción, corriente variacional, inversa de Holst,
% torsión algebraica, respuesta cuadrática y normalización métrica.
% Fuentes leídas íntegramente:
% [II10] output/ARTICULO_II_REV10_EDICION_INTEGRADA_20260910/manuscrito/
%        sections/10_gravedad.tex, líneas 109--155.
% [G11c] integral activo 2249, fuente/manuscrito/sections/md/
%        11c_gravedad_holonomia_rev6.tex, líneas 214--334.
% [B028] integral activo, fuente/manuscrito/integracion_83/parte_iv/body/
%        B028_ch68_realizacion_cartan_holst_barbero_7df2afd9f675_11_gravedad_cartan_holst.tex,
%        líneas 76--218; copia idéntica conservada en II10/fuentes/propietarios_II/G03/.
%        SHA-256 e1e1f3325399b3ef09c331ed8268e3883f7bcc561cd7a560cf80d1be7eb1e91c.
% CERTIFICADO_NUEVO de resultados recuperados: inversión explícita del mapa
% torsión -> D(e wedge e), prueba de unicidad de contorsión y eliminación
% homogénea del término cuadrático. Las dos inversas se comprobaron sobre
% las 24 bases correspondientes mediante aritmética racional exacta.
% Precisión de dominio: la respuesta lineal en corriente y su eliminación
% cuadrática se enuncian para materia afín en la contorsión. Se conserva
% la dependencia del coeficiente respecto de la acción material concreta.
%
% DEPENDENCIAS DE INTEGRACIÓN (no constituyen remisiones bibliográficas):
% 1. vii:sec:realizacion-cartan: cotetrada, conexión y realización del
%    transporte enriquecido; incluir su prueba en el cuerpo precedente.
% 2. vii:sec:acoplamientos: genealogía de gamma_HMT, G_int, c_int,
%    homogeneidad dimensional y elección de la carta, con pruebas.
% 3. Para evaluar una especie material: acción S_m concreta, representación
%    de sus campos y cálculo de su corriente sigma. Este fragmento recibe
%    esa acción declarada y demuestra su respuesta variacional.
% 4. La identificación de una corriente de pantalla con sigma, su
%    contracción física y su normalización material tienen residencia propia.
%    Aquí no se introduce ni se selecciona una amplitud s_0.
% 5. El balance métrico utiliza la covariancia de la acción material y
%    sus ecuaciones de movimiento; no supone conservación separada del
%    tensor de espín y del tensor material de Levi--Civita.

\section{Spin current and algebraic Cartan--Holst response}
\label{vii:vii:sec:corriente-respuesta-cartan}

The realization of the enriched transport in Section~\ref{vii:vii:sec:realizacion-cartan} provides the cotetrad and the connection; the dimensional sections and the incidence parameter of Sections~\ref{vii:v:ant:sec:barbero} and~\ref{vii:v:ant:sec:gravity} fix their couplings. Their composition determines a first-order variational problem. The material current occupies a precise position in it: it is the derivative of the action with respect to the connection and determines the torsional response with the same normalization.

\subsection{Geometric domain and current normalization}
\label{vii:vii:subsec:dominio-corriente-cartan}

We work on a smooth oriented four-dimensional manifold \(M\), with internal Lorentzian bundle \(E\), metric \(\eta=\operatorname{diag}(-1,1,1,1)\) and non-degenerate cotetrad \(e\in\Omega^1(M;E)\). In a local frame we write \(e^I\), with \(I=0,1,2,3\). The metric connection \(\omega^{IJ}=-\omega^{JI}\) acts through \(D_\omega\). When the fields \(\Psi\) are spinorial, a spin structure and the corresponding material representation are preserved. We define
\begin{equation}
 \Sigma^{IJ}=e^I\wedge e^J,\qquad
 F^{IJ}=\mathrm d\omega^{IJ}+\omega^I{}_{K}\wedge\omega^{KJ},
 \qquad T^I=D_\omega e^I.
 \label{vii:vii:eq:campos-cartan}
\end{equation}
The operator \(\star\) acts on the internal bivector indices,
\begin{equation}
 (\star X)^{IJ}=\tfrac12\epsilon^{IJ}{}_{KL}X^{KL},
 \qquad \star^2=-I,\qquad
 \mathsf P_\gamma=\star+\gamma^{-1}I.
 \label{vii:vii:eq:operador-holst}
\end{equation}
This distinguishes this duality from the exterior duality of forms on \(M\). Internal indices are raised and contracted with \(\eta\).

We fix a chart with \(G_{\rm int}>0\), \(c_{\rm int}>0\),
\(\gamma=\gamma_{\rm HMT}\in\mathbb R\setminus\{0\}\) and
\(\Lambda_{\rm int}\) constant on it. We abbreviate
\(\kappa=8\pi G_{\rm int}/c_{\rm int}^4\) and \(c=c_{\rm int}\).
These sections remain fixed during variation. The action is
\begin{equation}
 \begin{split}
 S[e,\omega,\Psi]={}&
 \frac1{2\kappa c}\int_M\Sigma_{IJ}\wedge
                         (\mathsf P_\gamma F)^{IJ}
 -\frac{\Lambda_{\rm int}}{\kappa c}
                         \int_M\operatorname{vol}_e
 +S_{\rm m}[e,\omega,\Psi].
 \end{split}
 \label{vii:vii:eq:accion-cartan-holst}
\end{equation}
The volume and orientation convention is the one that gives \(\Sigma_{IJ}\wedge(\star F)^{IJ}=R\operatorname{vol}_e\) in the Levi--Civita sector. The prior dimensional type makes each summand an action section.

For the variation of the connection, variations with compact support in the interior are used, or \(\delta\omega|_{\partial M}=0\) is fixed. A boundary completion whose variational problem cancels the same boundary term may also be used. At fixed cotetrad and material fields, the spin current is the degree-three form, valued in dual internal bivectors, determined by
\begin{equation}
 \delta_\omega S_{\rm m}
 =-\tfrac12\int_M\delta\omega^{IJ}\wedge\sigma_{IJ}.
 \label{vii:vii:eq:corriente-variacional}
\end{equation}
Pairing with an arbitrary compactly supported variation fixes \(\sigma\) uniquely. This definition preserves the action unit of the current and the normalization factor involved in its geometric response.

\subsection{Connection variation and Holst inversion}
\label{vii:vii:subsec:variacion-holst}

\begin{theorem}[Variational Cartan--Holst equation]
\label{vii:vii:thm:ecuacion-cartan-holst}
On the preceding domain, stationarity with respect to the connection is equivalent to
\begin{equation}
 D_\omega(\mathsf P_\gamma\Sigma)^{IJ}
       =\kappa c\,\sigma^{IJ}.
 \label{vii:vii:eq:ecuacion-cartan-holst}
\end{equation}
\end{theorem}
\begin{proof}
The variation of the curvature is \(\delta F=D_\omega\delta\omega\). The internal duality is parallel and self-adjoint for the bivector contraction; since \(\gamma\) is constant, \(D_\omega\mathsf P_\gamma=\mathsf P_\gamma D_\omega\). Setting \(B_{IJ}=(\mathsf P_\gamma\Sigma)_{IJ}\) yields
\[
 B_{IJ}\wedge D_\omega\delta\omega^{IJ}
 =\mathrm d(B_{IJ}\wedge\delta\omega^{IJ})
   -D_\omega B_{IJ}\wedge\delta\omega^{IJ}.
\]
The last form changes sign when its factors of degrees three and one are interchanged. The boundary conditions eliminate the integral of the exact term. The cosmological term is independent of \(\omega\). Therefore,
\[
 \delta_\omega S=
 \frac1{2\kappa c}\int_M\delta\omega^{IJ}\wedge
       \bigl[D_\omega B_{IJ}-\kappa c\,\sigma_{IJ}\bigr].
\]
The arbitrariness of the variation proves the equation, including its
sign and coefficient.
\end{proof}

\begin{lemma}[Real inverse of the Holst operator]
\label{vii:vii:lem:inversa-holst}
For \(\gamma\in\mathbb R\setminus\{0\}\),
\begin{equation}
 \mathsf P_\gamma^{-1}
 =\frac{\gamma^2}{1+\gamma^2}
                      (-\star+\gamma^{-1}I).
 \label{vii:vii:eq:inversa-holst}
\end{equation}
\end{lemma}
\begin{proof}
Both factors are polynomials in \(\star\) and commute. Their
product before normalization is
\[
 (\star+\gamma^{-1}I)(-\star+\gamma^{-1}I)
 =-\star^2+\gamma^{-2}I=(1+\gamma^{-2})I.
\]
The denominator \(1+\gamma^2\) is positive on the fixed real domain. Multiplication by \(\gamma^2/(1+\gamma^2)\) gives the identity.
\end{proof}

The current thus determines the degree-three bivector-valued form
\begin{equation}
 Y^{IJ}:=\kappa c\,
   \frac{\gamma^2}{1+\gamma^2}
   (-\star+\gamma^{-1}I)\sigma^{IJ},
 \qquad D_\omega\Sigma^{IJ}=Y^{IJ}.
 \label{vii:vii:eq:corriente-dualizada}
\end{equation}

\subsection{Pointwise reconstruction of torsion and contorsion}
\label{vii:vii:subsec:reconstruccion-torsion}

The Leibniz rule gives
\begin{equation}
 D_\omega\Sigma^{IJ}
 =T^I\wedge e^J-e^I\wedge T^J
 =:\mathsf L_e(T)^{IJ}.
 \label{vii:vii:eq:mapa-torsion}
\end{equation}
At each point, \(\mathsf L_e\) acts from \(\Lambda^2T^*M\otimes E\) to \(\Lambda^3T^*M\otimes\Lambda^2E\). Both spaces have dimension twenty-four. Let \(E_I\) be the dual frame of the cotetrad, so that \(\iota_{E_I}e^J=\delta_I^J\).

\begin{theorem}[Algebraic inversion of the torsional equation]
\label{vii:vii:thm:inversion-torsion}
For a nondegenerate co-tetrad, \(\mathsf L_e\) is an isomorphism.
Its inverse on the form \(Y\) is written
\begin{equation}
 B^I=\sum_J\iota_{E_J}Y^{IJ},\qquad
 b=\tfrac14\sum_I\iota_{E_I}B^I,\qquad
 T^I=B^I-e^I\wedge b.
 \label{vii:vii:eq:torsion-explicita}
\end{equation}
Equation~\eqref{vii:vii:eq:ecuacion-cartan-holst} therefore determines
torsion pointwise from \(e\) and the current appearing in it.
\end{theorem}
\begin{proof}
First let \(Y=\mathsf L_e(T)\) and set \(t=\sum_J\iota_{E_J}T^J\). The contraction of \eqref{vii:vii:eq:mapa-torsion}, using that \(T^I\) has degree two and that the cotetrad has four elements, gives
\[
 \sum_J\iota_{E_J}(T^I\wedge e^J)=2T^I,
 \qquad
 \sum_J\iota_{E_J}(e^I\wedge T^J)=T^I-e^I\wedge t.
\]
Hence \(B^I=T^I+e^I\wedge t\). A second contraction produces
\[
 \sum_I\iota_{E_I}B^I=t+3t=4t.
\]
One recovers \(t=b\) and then \(T^I=B^I-e^I\wedge b\). In particular, the kernel of \(\mathsf L_e\) is zero. Equality of dimensions proves its surjectivity and makes the formula valid for every form \(Y\) in the codomain.
\end{proof}

The connection decomposes uniquely as
\begin{equation}
 \omega^{IJ}=\mathring\omega^{IJ}(e)+K^{IJ},\qquad
 K^{IJ}=-K^{JI},\qquad T^I=K^I{}_{J}\wedge e^J,
 \label{vii:vii:eq:contorsion}
\end{equation}
where \(\mathring\omega(e)\) is the Levi--Civita connection and \(K\) the contorsion. To make the inverse explicit, we write \(K_{IJ}=K_{IJK}e^K\) and \(T_I=\tfrac12T_{IJK}e^J\wedge e^K\). Then
\begin{equation}
 T_{IJK}=K_{IKJ}-K_{IJK},\qquad
 K_{IJK}=\tfrac12(T_{KIJ}-T_{IJK}-T_{JKI}).
 \label{vii:vii:eq:contorsion-componentes}
\end{equation}
The second equality is obtained by combining the first cyclically and using \(K_{IJK}=-K_{JIK}\); substituting it recovers the components of \(T\). Torsion and contorsion are therefore two equivalent presentations of the algebraic response at fixed cotetrad.

\begin{corollary}[Current-free sector and minimal linear response]
\label{vii:vii:cor:desacoplamiento-torsional}
If \(\sigma=0\), one obtains \(T=0\) and \(\omega=\mathring\omega(e)\). For a material coupling whose current \(\sigma[e,\Psi]\) is independent of \(K\), equations~\eqref{vii:vii:eq:corriente-dualizada}, \eqref{vii:vii:eq:torsion-explicita} and \eqref{vii:vii:eq:contorsion-componentes} provide a unique response \(K_*(e,\Psi)\), linear in \(\sigma\).
\end{corollary}
\begin{proof}
The three reconstruction maps are linear with \(e\), \(\gamma\) and \(\kappa c\) fixed. The zero current produces \(Y=0\), then \(T=0\) and finally \(K=0\).
\end{proof}

In the minimal action, the connection equation contains this pointwise response. Its evolution is obtained by evolving \(e\) and \(\Psi\): \(K_*=K_*(e,\Psi)\) changes with them. If a material action preserves additional algebraic dependence on \(K\), the Cartan equation retains its variational form and must be solved with that current; the preceding linear response corresponds to the specified minimal domain.

\subsection{Elimination of contorsion and quadratic contribution}
\label{vii:vii:subsec:respuesta-cuadratica}

Now consider the minimal coupling affine in \(K\), expressed by
\begin{equation}
 S_{\rm m}[e,\mathring\omega+K,\Psi]
 =S_{\rm m}[e,\mathring\omega,\Psi]
   -\tfrac12\int_M K^{IJ}\wedge\sigma_{IJ}[e,\Psi].
 \label{vii:vii:eq:materia-afin-contorsion}
\end{equation}
This condition identifies the material action to which the quadratic elimination applies; the current remains the one defined by \eqref{vii:vii:eq:corriente-variacional}.

The curvature expansion is
\begin{equation}
 F(\mathring\omega+K)
 =F(\mathring\omega)+D_{\mathring\omega}K+K\wedge K,
 \qquad (K\wedge K)^{IJ}=K^I{}_{L}\wedge K^{LJ}.
 \label{vii:vii:eq:curvatura-contorsion}
\end{equation}
Since \(D_{\mathring\omega}\Sigma=0\), the linear term is
the exterior derivative of \(\Sigma_{IJ}\wedge(\mathsf P_\gamma K)^{IJ}\).
The corresponding boundary contribution remains explicit:
\begin{equation}
 S_\partial[e,K]=\frac1{2\kappa c}
     \int_{\partial M}\Sigma_{IJ}\wedge(\mathsf P_\gamma K)^{IJ}.
 \label{vii:vii:eq:borde-contorsion}
\end{equation}
On an interior domain, or with the compatible boundary completion, the action density that depends on the contorsion is
\begin{equation}
 \mathcal Q_{e,\gamma}(K)-\tfrac12K^{IJ}\wedge\sigma_{IJ},
 \qquad
 \mathcal Q_{e,\gamma}(K)
  =\frac1{2\kappa c}\Sigma_{IJ}\wedge
                         \mathsf P_\gamma(K\wedge K)^{IJ}.
 \label{vii:vii:eq:forma-cuadratica-contorsion}
\end{equation}
The density \(\mathcal Q_{e,\gamma}\) is a degree-four differential form, homogeneous of degree two in \(K\).

\begin{theorem}[Effective response quadratic in the current]
\label{vii:vii:thm:respuesta-cuadratica}
Under \eqref{vii:vii:eq:materia-afin-contorsion}, substituting the algebraic solution \(K_*\) produces the effective spin density
\begin{equation}
 \mathcal L_{\rm spin}^{\rm eff}
 =-\tfrac14 K_*^{IJ}\wedge\sigma_{IJ}
 =-\mathcal Q_{e,\gamma}(K_*).
 \label{vii:vii:eq:accion-efectiva-spin}
\end{equation}
At fixed cotetrad and couplings, it is homogeneous of degree two in the current. The effective action also preserves \(S_\partial[e,K_*]\), or its declared boundary completion.
\end{theorem}
\begin{proof}
By Corollary~\ref{vii:vii:cor:desacoplamiento-torsional}, \(K_*\) is linear in \(\sigma\). The connection equation is precisely the stationarity of \eqref{vii:vii:eq:forma-cuadratica-contorsion}. Since it is algebraic, the radial variation \(\delta K=K_*\) is applied pointwise; it may also be multiplied by a compactly supported function. Euler's identity for a quadratic form gives
\[
 2\mathcal Q_{e,\gamma}(K_*)
             -\tfrac12K_*^{IJ}\wedge\sigma_{IJ}=0.
\]
Substitution into the action density proves \eqref{vii:vii:eq:accion-efectiva-spin}. If \(\sigma\mapsto\lambda\sigma\), then \(K_*\mapsto\lambda K_*\) and the effective density is multiplied by \(\lambda^2\). The expansion \eqref{vii:vii:eq:curvatura-contorsion} separately identifies the boundary term accompanying this elimination.
\end{proof}

The specific current, its Lorentzian contraction, the duality and the conventions of the material action determine the sign and coefficient of this contribution. The quadratic identity provides the composition rule that allows them to be evaluated once these data have been specified.

\subsection{Effective metric source and covariant balance}
\label{vii:vii:subsec:fuente-metrica-efectiva}

The normalization of the metric source is fixed in the same action line. After eliminating \(K\), we define
\begin{align}
 \delta_g S_{\rm m}[e,\mathring\omega,\Psi]
 &=-\tfrac12\int_M\operatorname{vol}_e\,
                   \mathcal T^{S,\rm LC}_{\mu\nu}\,\delta g^{\mu\nu},
 \\
 \delta_g S_{\rm spin}^{\rm eff}
 &=-\tfrac12\int_M\operatorname{vol}_e\,
                   \mathcal U^{S,\rm spin}_{\mu\nu}\,\delta g^{\mu\nu}.
 \label{vii:vii:eq:fuentes-normalizadas-accion}
\end{align}
Here interior variations are considered, with boundary terms treated according to the fixed variational problem. For the affine coupling, the second source is the metric variation of \eqref{vii:vii:eq:accion-efectiva-spin}. The variation includes the dependence of the material current on \(e\) and \(\Psi\).

\begin{proposition}[Metric equation and conservation of the total source]
\label{vii:vii:prop:fuente-metrica-balance}
With the preceding conventions, the metric equations are
\begin{equation}
 G_{\mu\nu}[g]+\Lambda_{\rm int}g_{\mu\nu}
 =\kappa c\bigl(\mathcal T^{S,\rm LC}_{\mu\nu}
                    +\mathcal U^{S,\rm spin}_{\mu\nu}\bigr)
 =\kappa\bigl(T^{\rm phys,LC}_{\mu\nu}
                    +U^{\rm phys,spin}_{\mu\nu}\bigr),
 \label{vii:vii:eq:ecuacion-metrica-efectiva}
\end{equation}
where \(T^{\rm phys,LC}=c\mathcal T^{S,\rm LC}\) and \(U^{\rm phys,spin}=c\mathcal U^{S,\rm spin}\) when the cotetrad takes values in the length line. Every solution satisfies
\begin{equation}
 \mathring\nabla^\mu
   (T^{\rm phys,LC}_{\mu\nu}+U^{\rm phys,spin}_{\mu\nu})=0.
 \label{vii:vii:eq:balance-metrico-total}
\end{equation}
\end{proposition}
\begin{proof}
Eliminating \(K\) preserves the equations of the remaining variables: in the chain rule, the term multiplying \(\delta K_*\) vanishes by the connection equation. In the Levi--Civita sector, the first Bianchi identity cancels the term \(\Sigma_{IJ}\wedge F^{IJ}\), and the gravitational part is written \((2\kappa c)^{-1}\int_M(R-2\Lambda_{\rm int})\operatorname{vol}_e\). Its interior variation is
\[
 \frac1{2\kappa c}\int_M\operatorname{vol}_e\,
        (G_{\mu\nu}+\Lambda_{\rm int}g_{\mu\nu})\delta g^{\mu\nu}.
\]
The sum with \eqref{vii:vii:eq:fuentes-normalizadas-accion} gives
\eqref{vii:vii:eq:ecuacion-metrica-efectiva}. The contracted Bianchi identity
for \(\mathring\nabla\), metric compatibility and
the constancy of \(\kappa\) and \(\Lambda_{\rm int}\) in the chart
produce \eqref{vii:vii:eq:balance-metrico-total}.
\end{proof}

Before eliminating the connection, the geometric identities retain
the form
\begin{equation}
 D_\omega T^I=F^I{}_{J}\wedge e^J,\qquad D_\omega F^{IJ}=0.
 \label{vii:vii:eq:bianchi-cartan}
\end{equation}
The first is \(D_\omega^2e=F\wedge e\); the second is obtained by expanding \(D_\omega(\mathrm d\omega+\omega\wedge\omega)\) and cancelling the terms. In a covariant material action, the total metric balance coincides with the Noether identity on the material equations. Geometry and matter thus preserve a joint source; the partition between its Levi--Civita contribution and its spin contribution is determined by the specialized action.
